% Generated draft: requires mathematical and source review.
\documentclass[11pt]{article}
\usepackage[margin=25mm]{geometry}
\usepackage{fontspec}
\setmainfont{Latin Modern Roman}
\usepackage{amsmath,amssymb,mathtools,unicode-math}
\setmathfont{Latin Modern Math}
\usepackage{enumitem,needspace,xurl,hyperref}
\setlength{\parindent}{0pt}
\setlength{\parskip}{5pt}
\setlength{\emergencystretch}{3em}
\newcommand{\N}{\mathbb{N}}
\newcommand{\Z}{\mathbb{Z}}
\newcommand{\Q}{\mathbb{Q}}
\newcommand{\R}{\mathbb{R}}
\newcommand{\C}{\mathbb{C}}
\newcommand{\F}{\mathbb{F}}
\newcommand{\A}{\mathbb{A}}
\newcommand{\PP}{\mathbb{P}}
\newcommand{\E}{\mathbb{E}}
\newcommand{\eps}{\varepsilon}
\newcommand{\dd}{\,\mathrm{d}}
\begin{document}
\section*{UnsolvedMath ID: 20000410}
\textbf{Source problem code:} catalogue entry
\subsection{Definitions and mathematical statement}
The source's comment asks to make a surface over \(\F_q\) into a curve over \(\F_q(T)\) by a Lefschetz pencil and to ``see'' an expected second Tate class beyond the hyperplane class, or the corresponding Jacobian point. For a smooth projective surface, a Tate class is understood in \(H^2_{\mathrm{et}}(S_{\overline{\F}_q},\Q_\ell(1))\), \(\ell\nmid q\), fixed by Frobenius or by a specified power of it. Those two field-of-definition conventions differ. The source's ``typically'' two-class/parity remark is not an assumption that every surface has even second Betti number or exactly two classes; an arbitrary extra algebraic class is not supplied.

For the attached aid's conditional construction, take a smooth projective surface \(S/\F_q\), an ample pencil class \(H\), and a blowup \(\pi:\widetilde S\to S\) resolving a Lefschetz pencil \(\widetilde S\to\PP^1\), with smooth generic curve \(C/K\), \(K=\F_q(T)\), and Jacobian \(J_C\). A divisor class \(D_0\) is \(H\)-primitive if \(D_0\cdot H=0\), so restriction gives a degree-zero class on \(C\). The aid reports that this point is torsion exactly when \(\pi^*D_0\) lies in the rational span of geometric fiber-component classes. In particular, when all geometric fibers are integral, every nonzero primitive N\'eron--Severi class restricts to a non-torsion point, under the pencil hypotheses. If a divisor is defined only after a finite constant-field extension, its resulting point is over that extension of \(K\), not automatically over \(K\).

For a divisor \(D\) defined over the chosen ground field, put
\[
D_0=(H^2)D-(D\cdot H)H.
\]
The aid's numerical certificate is
\[
H^2D^2-(D\cdot H)^2<0
\quad\Longrightarrow\quad
[\mathcal O_C(D_0)]\text{ non-torsion},
\]
assuming the integral-geometric-fiber condition. This certifies a point from an already available divisor; it does not construct an unknown Tate class on every surface.

Its explicit specialization is a smooth quartic K3 surface \(S\subset\PP^3\) containing a line \(L\), with a suitable plane pencil of axis \(R\) skew to \(L\). Set \(p_t=L\cap\Pi_t\) and \(B=S\cap R\), viewed on the generic plane quartic. The aid reports
\[
P=[4p_t-B]\in J_C(K)
\]
as non-torsion when its pencil hypotheses hold. Here \(\deg B=4\), \(H^2=4\), \(L\cdot H=1\), \(L^2=-2\), and the certificate is \(-9<0\). Ground-field definition of the line/pencil and geometric integrality of every fiber remain required. The source's intended general-surface problem remains unresolved by this conditional example.
\subsection{Short English statement}
Construct the expected extra class or Jacobian point arising from a surface Lefschetz pencil. The attached aid reports a torsion criterion for restrictions of primitive divisor classes and a conditional explicit non-torsion point on a quartic K3 pencil, without providing a new class for every surface.
\subsection{Sources}
\begin{itemize}
\item Local UnsolvedMath registry, permanent ID \(20000410\), exact public code \texttt{catalogue entry}; canonical location \texttt{aim-algebraic-number-theory-notes.json}, notes[61].
\item ``Problems from the AIM Tate Conjecture Workshop,'' July 23--27, 2007, transcribed by Christopher Lyons; supplied URL \url{https://aimath.org/WWN/tateconjecture/tateconjecture.pdf}. The verification identifies this as the second comment on C. Hall's Question 1, attributed to C. Schoen, extending the preceding hyperelliptic-Jacobian question. The PDF was not consulted here; the original literature field is empty.
\item Attached research attempt 1, ``Seeing a primitive surface class in the generic Jacobian,'' without full bibliographic references for the restriction criterion or pencil construction.
\end{itemize}
\end{document}
